\chapter{Ղեկավարում, համագործակցություն և ռիսկեր}
\label{ch:supervision}

\section{Ղեկավարում}
\begin{itemize}
  \item \textbf{\supervisorlabel{1}} (Միկրոէլեկտրոնիկա) — RTL, FPGA, վերիֆիկացիա։
  \item \textbf{\supervisorlabel{2}} (Ծրագրային ճարտարագիտություն) — Backend, տվյալներ, վերլուծություն։
  \item Ընդհանուր հանդիպում՝ շաբաթը մեկ անգամ, ինտեգրման փուլում՝ երկու անգամ։
\end{itemize}

\section{Ռիսկեր}
\begin{risks}[Հիմնական ռիսկեր և կանխարգելում.]
  \risk{AXI/Zynq ինտեգրման բարդությունը}{բարձր}{Ներածական վարժություն 2, պատրաստի AXI IP շաբլոններ, ղեկավարի աջակցություն}
  \risk{Ուսանող 2-ը սպասում է ապարատին}{միջին}{Էմուլյատոր-սկրիպտ նույն JSON ֆորմատով (Խ2.3)}
  \risk{Timing-ը չի անցնում}{միջին}{Pipeline, ցածր տակտային հաճախություն MVP-ի համար}
  \risk{Հարթակը հասանելի չէ}{ցածր}{Սիմուլյացիա, ժամանակացույց լաբորատորիայում}
\end{risks}
